\documentclass[11pt]{article}
\usepackage[T1]{fontenc}
\usepackage[margin=1in]{geometry}
\usepackage{amsmath}
\usepackage{booktabs}
\usepackage{siunitx}

\title{Results: Thermal Conductivity Measurements Across Sample Batches}
\author{Materials Characterization Group}
\date{}

\begin{document}
\maketitle

\section{Introduction}
This document reports thermal conductivity, density, and tensile strength
measurements collected across four sample batches of a composite laminate
material. All measurements were performed under controlled laboratory
conditions, with uncertainties reported at the one-standard-deviation
level unless otherwise noted. Section~\ref{sec:conductivity} reports the
conductivity results using the modern \texttt{siunitx} unit macros, while
Section~\ref{sec:density} reports density and strength results using the
legacy \texttt{\textbackslash SI} and \texttt{\textbackslash si} macros for
comparison, since both forms remain common in the literature this study
draws on.

\section{Thermal Conductivity}
\label{sec:conductivity}
Table~\ref{tab:conductivity} summarizes the measured thermal conductivity
for each batch at three test temperatures, given as
\qty{25}{\celsius}, \qty{75}{\celsius}, and \qty{125}{\celsius}. Values are
reported in \unit{\watt\per\meter\per\kelvin}, with sample counts and
standard deviations given alongside each mean.

\begin{table}[h]
\centering
\caption{Thermal conductivity (\unit{\watt\per\meter\per\kelvin}) by batch and
temperature.}
\label{tab:conductivity}
\begin{tabular}{l S[table-format=1.3] S[table-format=1.3] S[table-format=1.3]}
\toprule
{Batch} & {\qty{25}{\celsius}} & {\qty{75}{\celsius}} & {\qty{125}{\celsius}} \\
\midrule
A1 & 0.842 & 0.901 & 0.958 \\
A2 & 0.831 & 0.889 & 0.947 \\
B1 & 0.917 & 0.972 & 1.034 \\
B2 & 0.905 & 0.963 & 1.021 \\
\bottomrule
\end{tabular}
\end{table}

Across all four batches, conductivity increases approximately linearly
with temperature over the tested range, and batches B1 and B2 consistently
outperform A1 and A2, which is consistent with the higher filler fraction
used in the B-series formulation. The largest single measurement,
\qty{1.034}{\watt\per\meter\per\kelvin} for batch B1 at
\qty{125}{\celsius}, remains well within the manufacturer's rated
operating envelope of \qty{1.2}{\watt\per\meter\per\kelvin}.

Table~\ref{tab:stats} reports summary statistics across all batches
combined, including the sample count, mean, and standard deviation at each
temperature, using \texttt{multicolumn} and \texttt{cmidrule} to group the
per-temperature statistics under shared headers.

\begin{table}[h]
\centering
\caption{Summary statistics across all batches.}
\label{tab:stats}
\begin{tabular}{l ccc ccc ccc}
\toprule
& \multicolumn{3}{c}{\qty{25}{\celsius}} & \multicolumn{3}{c}{\qty{75}{\celsius}} & \multicolumn{3}{c}{\qty{125}{\celsius}} \\
\cmidrule(lr){2-4} \cmidrule(lr){5-7} \cmidrule(lr){8-10}
{} & {$n$} & {Mean} & {SD} & {$n$} & {Mean} & {SD} & {$n$} & {Mean} & {SD} \\
\midrule
All batches & 4 & 0.874 & 0.041 & 4 & 0.931 & 0.040 & 4 & 0.990 & 0.041 \\
\bottomrule
\end{tabular}
\end{table}

The standard deviation at each temperature is small relative to the mean,
of order \num{4e-2}, indicating good repeatability across the four
batches. We note that the sample size per batch was $n = \num{12}$
individual coupons, so the batch means reported in
Table~\ref{tab:conductivity} are themselves averages over a reasonably
large number of measurements, and the pooled standard deviations in
Table~\ref{tab:stats} should be interpreted accordingly.

\section{Density and Tensile Strength}
\label{sec:density}
For comparison with legacy laboratory notebooks, this section reports
density and tensile strength using the older \si{} and \SI{} macro forms.
Density was measured as \SI{1.58}{\gram\per\centi\meter\cubed} for the
A-series batches and \SI{1.64}{\gram\per\centi\meter\cubed} for the
B-series batches, consistent with the higher filler loading noted above.
Table~\ref{tab:strength} reports tensile strength alongside elongation at
break for all four batches.

\begin{table}[h]
\centering
\caption{Tensile strength and elongation at break.}
\label{tab:strength}
\begin{tabular}{l S[table-format=3.1] S[table-format=1.2]}
\toprule
{Batch} & {Strength (\si{\mega\pascal})} & {Elongation (\si{\percent})} \\
\midrule
A1 & 142.3 & 2.11 \\
A2 & 138.7 & 2.04 \\
B1 & 156.9 & 1.87 \\
B2 & 151.2 & 1.92 \\
\bottomrule
\end{tabular}
\end{table}

The B-series batches show higher tensile strength but lower elongation at
break than the A-series batches, a trade-off that is typical of increased
filler loading in composite laminates: the additional filler stiffens the
matrix and raises peak strength, at the cost of the ductility that allows
the A-series material to elongate further before fracture. Combined with
the conductivity results in Section~\ref{sec:conductivity}, these
measurements suggest that batch selection should be guided by the relative
importance of thermal performance versus mechanical ductility in the
intended application, since no single batch dominates on every axis
studied here.

\section{Conclusion}
Across all measured properties, the B-series formulation offers higher
thermal conductivity and tensile strength than the A-series formulation,
at the cost of reduced elongation at break. These results, summarized in
Tables~\ref{tab:conductivity}, \ref{tab:stats}, and \ref{tab:strength},
will inform the batch selection guidance issued alongside the next
revision of the material datasheet.

\end{document}
